\subsection{Terminal User Interface}
\label{app:tui}

The reference implementation ships both a command-line interface and a
full-screen terminal user interface.
The CLI covers one-shot scans, library inspection, and defense telemetry
without entering an interactive session.
The TUI exposes the same defense surface through slash commands, modal
overlays, and a streaming chat pane built on the same session model as a
modern coding agent.

\paragraph{CLI.}
Figure~\ref{fig:appendix-cli-help} lists the top-level command surface.
Figure~\ref{fig:appendix-cli-status} shows the antibody and antigen library
inventory.
Figure~\ref{fig:appendix-cli-dashboard} reports defense telemetry from the
CLI dashboard.
Figure~\ref{fig:appendix-cli-scan} records a single-shot scan that matches
Tier~0 signatures without an LLM call.
All four panels are drawn from live CLI transcripts on a development
workstation after the evaluation campaigns in
Section~\ref{sec:evaluation}.

\paragraph{TUI.}
Figure~\ref{fig:appendix-tui-cold} shows the cold-start screen with the
logo, status footer, and empty editor.
Figure~\ref{fig:appendix-tui-dashboard} shows the defense-dashboard overlay
opened with \texttt{/dashboard}.
Sessions persist as append-only JSONL files under
\texttt{\textasciitilde/.caitlyn/sessions/}, so every scan, vaccination, and
branch is reproducible from the session log alone.
The dashboard reads the same statistics events that drive System~II
triggers, which lets operators observe anomaly signals before an immune
response is raised.
Table~\ref{tab:appendix-tui-commands} lists the defense-facing and
session slash commands.

\begin{table}[h]
  \centering
  \footnotesize
  \setlength{\tabcolsep}{4pt}
  \renewcommand{\arraystretch}{1.18}
  \caption[Slash commands exposed by the CAITLYN terminal user interface.]{%
    Slash commands exposed by the CAITLYN terminal user interface.
    The upper block lists defense-facing commands.
    The lower block lists session and configuration commands.
  }
  \label{tab:appendix-tui-commands}
  \rowcolors{2}{BoxGray!30}{white}
  \begin{tabular*}{\columnwidth}{@{\extracolsep{\fill}}
    >{\ttfamily\bfseries\raggedright\arraybackslash}p{0.46\columnwidth}
    >{\raggedright\arraybackslash}p{0.42\columnwidth}
    @{}}
    \toprule
    \multicolumn{1}{@{}l}{\textbf{Command}} &
    \multicolumn{1}{l@{}}{\textbf{Purpose}} \\
    \midrule
    \rowcolor{white}
    \multicolumn{2}{@{}l@{}}{\emph{Defense}} \\
    /scan <content> & Single-shot scan \\
    /status & Antibody and antigen library \\
    /dashboard & Defense statistics \\
    /history [N] & Recent scan history \\
    /antibody list|add|remove & Library edits \\
    /antigen <id> & Antigen details \\
    /vaccinate <pattern> & Immune response \\
    \midrule
    \rowcolor{white}
    \multicolumn{2}{@{}l@{}}{\emph{Session}} \\
    /new, /resume, /session & Session management \\
    /fork, /truncate, /tree & Branch navigation \\
    /model, /thinking & Model controls \\
    /login <provider> & Credential configuration \\
    /settings, /help & Configuration and help \\
    /export, /compact & Session artifacts \\
    /clear, /quit & Session control \\
    \bottomrule
  \end{tabular*}
\end{table}

